Write $A = A_{k}$ and $N = N_{k}$ when the rung is fixed. An element is \emph{even} if it is a combination of blades with an even number of factors.

\emph{Two facts about the algebras $A_{I}$.} (F1) For disjoint $I$, $J$ the map $A_{I} \otimes A_{J} \to A_{I \cup J}$, $a \otimes b \mapsto ab$, is a linear bijection: a blade of $A_{I \cup J}$ is a product of distinct generators, and moving those of index in $I$ to the front changes it by a sign only, so the products of a blade of $A_{I}$ and a blade of $A_{J}$ are, up to sign, the blades of $A_{I \cup J}$, each once (the blades are a basis by Theorem \claimref{synthesis.signed-directions-and-the-integer-clifford-algebra}(vii)). (F2) A generator of $A_{I}$ anticommutes with each generator of $A_{J}$, so it commutes with every even element of $A_{J}$.

\emph{(i), the isomorphism.} By induction on $k$. For $k = 1$ let $M$ be the bit over $\Q$ with set $v$ and clear $w$. By Theorem \claimref{synthesis.set-and-clear-of-a-bit-generate-the-hyperbolic-plane}(iv) the maps $v + w$ and $v - w$ anticommute and square to $1$ and $-1$, so the universal property of $\Cl_{\Q}(1,1)$ gives a homomorphism $A_{1} \to \mathrm{End}_{\Q}(M)$ with $e_{1} \mapsto v + w$, $f_{1} \mapsto v - w$. Its image contains $v$ and $w$ (here $\tfrac12$ is used), which generate $\mathrm{End}_{\Q}(M)$ by clause (ii) of the same theorem, and both algebras have dimension $4$; so it is an isomorphism. For $k \ge 2$ put $\omega = e_{k}f_{k}$. Then $\omega^{2} = 1$, $\omega$ anticommutes with $e_{k}$ and $f_{k}$ and commutes with every generator $x$ of $A_{k-1}$. So the elements $x\omega$ satisfy $(x\omega)^{2} = x^{2}$ and $(x\omega)(x'\omega) = xx' = -(x'\omega)(x\omega)$ for $x \ne x'$, and each commutes with $e_{k}$ and $f_{k}$. By the universal property there is a homomorphism $\phi : A_{k-1} \to A_{k}$ with $\phi(x) = x\omega$, whose image commutes with the subalgebra $C$ generated by $e_{k}, f_{k}$. Hence $a \otimes c \mapsto \phi(a)c$ is a homomorphism $A_{k-1} \otimes C \to A_{k}$; its image contains $e_{k}$, $f_{k}$, $\omega$ and each $x = (x\omega)\omega$, so it is onto, and both sides have dimension $4^{k}$. So $A_{k} \cong \mathrm{End}(S_{k-1}) \otimes \mathrm{End}(M) = \mathrm{End}(S_{k-1} \otimes M)$, and $S_{k} = S_{k-1} \otimes M$ has dimension $2^{k}$.

\emph{(i), the rank.} Identify $A$ with $\mathrm{End}(S_{k})$. Then $aA$ is the set of maps $t$ with $\operatorname{im} t \subseteq \operatorname{im} a$: one inclusion is clear, and if $\operatorname{im} t \subseteq \operatorname{im} a$ a map $s$ with $as = t$ is defined on a basis. That set has dimension $2^{k}\operatorname{rank}a$, and $\dim A = 4^{k}$. Every rank from $0$ to $2^{k}$ occurs.

\emph{(ii).} $aA$ is a subspace of $A$ containing $a$, which gives the bounds and $N(a) = 0$ only for $a = 0$. If $N(a) = 1$ then $ab = 1$ for some $b$; then $x \mapsto ax$ is onto, hence injective, and $a(ba - 1) = 0$ gives $ba = 1$. Conversely an invertible $a$ has $aA = A$. Next, $(a + b)A \subseteq aA + bA$ and $\lambda aA = aA$; $abA \subseteq aA$; and $abA$ is the image of $bA$ under $x \mapsto ax$, so its dimension is at most that of $bA$. For the last inequality let $Z = \{x : ax = 0\}$. The map $x \mapsto ax$ carries $bA$ onto $abA$ with kernel $bA \cap Z$ and $bcA$ onto $abcA$ with kernel $bcA \cap Z$, so $\dim bA - \dim abA = \dim(bA \cap Z)$ and $\dim bcA - \dim abcA = \dim(bcA \cap Z)$, and the second is at most the first because $bcA \subseteq bA$. Call these rules (N2) [the two for sums and multiples], (N3) [the two for products] and (N4) [the last].

\emph{Consequences for idempotents.} Let $M$ be any function with $M(0) = 0$, $M(1) = 1$, (N2) without the clause on multiples, (N3) and (N4). (a) If $e^{2} = e$ then $M(e) + M(1 - e) = 1$: $1 \le M(e) + M(1 - e)$ by (N2), and (N4) for $(e, 1, 1 - e)$ gives $M(e) + M(1 - e) \le M(1) + M(0)$. (b) If $e, f$ are idempotents with $ef = fe = 0$ then $M(e + f) = M(e) + M(f)$: one inequality is (N2), the other is (N4) for $(e, e + f, f)$, since $e(e + f) = e$, $(e + f)f = f$, $e(e + f)f = 0$. (c) If $u$ is invertible, $M(uau^{-1}) = M(a)$, by (N3) twice.

\emph{(iii) and (iv).} Let $I$, $J$ be disjoint, $a \in A_{I}$, and write $N_{I}$ for the size computed in $A_{I}$. By (F1), $aA_{I \cup J} = (aA_{I})A_{J}$ is the image of $(aA_{I}) \otimes A_{J}$ under a bijection, so $N_{I \cup J}(a) = N_{I}(a)$; with $I = \{1, \dots, k\}$, $J = \{k + 1\}$ this is (iii). If $b \in A_{J}$ is even it commutes with $A_{I}$ by (F2), so $ab = ba$ and $abA_{I \cup J} = abA_{I}A_{J} = (aA_{I})(bA_{J})$, the image of $(aA_{I}) \otimes (bA_{J})$ under the same bijection. Dividing its dimension by $\dim A_{I}\dim A_{J}$ gives $N(ab) = N(a)N(b)$.

\emph{(v).} Additivity is (b). For the tests, $e_{i}\,\mathrm{on}_{i}\,e_{i} = (e_{i}v_{i}e_{i})(e_{i}w_{i}e_{i}) = w_{i}v_{i} = \mathrm{off}_{i}$, because $e_{i}f_{i}e_{i} = -f_{i}$; so the two have the same size by (c), and the sizes add to $1$ by (a), since $\mathrm{on}_{i} + \mathrm{off}_{i} = v_{i}w_{i} + w_{i}v_{i} = 1$. For a null vector $n \ne 0$: $n^{2} = 0$ gives $nA \subseteq Z = \{x : nx = 0\}$, and $\dim nA + \dim Z = \dim A$, so $N(n) \le \tfrac12$. The form $Q$ is nondegenerate, so there is a vector $m$ with $nm + mn = 1$ after scaling. Then $1 \le N(nm) + N(mn) \le 2N(n)$.

\emph{(vi).} Fix a rung $A_{k} = \mathrm{End}(V)$, $\dim V = d = 2^{k}$. Two idempotents of the same rank are conjugate (choose bases adapted to image and kernel), so they have the same $M$ by (c). The identity is a sum of $d$ idempotents of rank $1$ with pairwise products zero, so by (b) each idempotent of rank $1$ has $M = 1/d$, and one of rank $r$ has $M = r/d$. For a general $a$ choose a basis $c_{1}, \dots, c_{r}$ of $\operatorname{im}a$, vectors $t_{i}$ with $at_{i} = c_{i}$, and a complement $D$ of $\operatorname{im}a$; let $y$ send $c_{i}$ to $t_{i}$ and vanish on $D$. Then $ay$ is the identity on $\operatorname{im}a$, so $aya = a$, and $e = ay$ is an idempotent of the rank of $a$ with $ea = a$. By (N3), $M(a) \le M(e) \le M(a)$. So $M(a) = \operatorname{rank}(a)/d = N(a)$.
